% Generated draft: requires mathematical and source review.
\documentclass[11pt]{article}
\usepackage[margin=25mm]{geometry}
\usepackage{fontspec}
\setmainfont{Latin Modern Roman}
\usepackage{amsmath,amssymb,mathtools,unicode-math}
\setmathfont{Latin Modern Math}
\usepackage{enumitem,needspace,xurl,hyperref}
\setlength{\parindent}{0pt}
\setlength{\parskip}{5pt}
\setlength{\emergencystretch}{3em}
\newcommand{\N}{\mathbb{N}}
\newcommand{\Z}{\mathbb{Z}}
\newcommand{\Q}{\mathbb{Q}}
\newcommand{\R}{\mathbb{R}}
\newcommand{\C}{\mathbb{C}}
\newcommand{\F}{\mathbb{F}}
\newcommand{\A}{\mathbb{A}}
\newcommand{\PP}{\mathbb{P}}
\newcommand{\E}{\mathbb{E}}
\newcommand{\eps}{\varepsilon}
\newcommand{\dd}{\,\mathrm{d}}
\begin{document}
\section*{UnsolvedMath ID: 10400185}
\textbf{Source problem code:} catalogue entry
\subsection{Definitions and mathematical statement}
Let \(\mathcal A(\varnothing;\Q)\) be the rational graded space of vertex-oriented trivalent Jacobi diagrams with no legs, modulo antisymmetry and IHX relations, with the usual degree equal to half the number of vertices. Write \(\mathcal A(\varnothing;\Q)^{(d)}_{\mathrm{conn}}\) for its degree-\(d\) connected part. The empty diagram has degree zero and belongs to the full space, not its connected part. A primitive invariant is additive under connected sum: \(v(M\mathbin{\#}N)=v(M)+v(N)\). Here \(M,N\) are integral homology \(3\)-spheres, and the finite-type and degree conventions are those intended by the source's diagrammatic equivalence.

\emph{Problem.} Determine, for every degree \(d\), the dimension of the degree-\(d\) space of primitive rational finite-type invariants of integral homology \(3\)-spheres; equivalently, determine
\[
 c_d=\dim_\Q\mathcal A(\varnothing;\Q)^{(d)}_{\mathrm{conn}}
 \quad\text{for every }d\geq0.
\]
The invariant formulation must use degree-\(d\) graded pieces, rather than treating the set of invariants of exactly degree \(d\) as a vector space. The omitted source setup is needed to fix its finite-type filtration and the matching of degree conventions.

Table 7 supplies the following values, where \(a_d\) denotes the dimension of the full diagram space and \(p_d\) retains the undefined table label ``prime diag.'' without interpreting it as another dimension:
\[
\begin{array}{c|rrrrrrrrrrr}
 d&0&1&2&3&4&5&6&7&8&9&10\\\hline
 p_d&0&1&0&0&1&0&1&1&1&1&2\\
 c_d&0&1&1&1&2&2&3&4&5&6&8\\
 a_d&1&1&2&3&6&9&16&25&42&65&105
\end{array}
\]
\[
\begin{array}{c|rrrr}
 d&11&12&13&14\\\hline
 c_d&9&\geq11&\geq13&\geq15\\
 a_d&161&\geq254&\geq386&\geq595
\end{array}
\]
The extracted continuation also has a ``prime diag.'' row containing a single \(1\), but does not preserve enough column alignment to assign that entry unambiguously among degrees \(11,12,13,14\). All inequalities above remain lower bounds, not exact values. The table is supplied historical evidence, not an independently verified present-day computation.
\subsection{Short English statement}
Compute the dimension in every degree of primitive finite-type invariants of integral homology \(3\)-spheres, equivalently of connected closed rational Jacobi diagrams. The supplied table gives small-degree values and higher-degree lower bounds, which are retained with their exact inequalities.
\subsection{Sources}
\begin{itemize}
\item Local UnsolvedMath registry, permanent ID 10400185, source code catalogue entry, Problem 10.10; no individual proposer is supplied.
\item Ohtsuki, \emph{Problems on invariants of knots and 3-manifolds} (2002), Problem 10.10 and Table 7, PDF page 146. Supplied URL: \url{https://msp.org/gtm/2002/04/gtm-2002-04-024s.pdf}; recorded access 2026-07-29; extraction pdf-font-structure. The collection distinguishes 231 numbered bold problem, conjecture, and question headings.
\item Supplied status partially\_solved and PARTIAL-PROGRESS; the summary says dimensions are known up to approximately degree \(10\), referring to Table 7. The actual supplied table also lists exact degree-\(11\) dimensions and bounds in degrees \(12\)--\(14\). Current primary-literature status and redistribution terms of the publicly accessible author/publisher copy remain NEEDS\_REVIEW.
\end{itemize}
\end{document}
